\documentclass{article}%
\usepackage[T1]{fontenc}%
\usepackage[utf8]{inputenc}%
\usepackage{lmodern}%
\usepackage{textcomp}%
\usepackage{lastpage}%
\usepackage{geometry}%
\geometry{margin=1in,headheight=14pt,headsep=25pt}%
\usepackage{hyperref}%
\usepackage{enumitem}%
\usepackage{fancyhdr}%
\usepackage{xcolor}%
\usepackage{url}%
\usepackage{breakurl}%
\usepackage{amsmath}%
\usepackage{amssymb}%
\usepackage{mathtools}%
\usepackage{amsthm}%
\usepackage{thmtools}%
\usepackage{float}%
\usepackage{subcaption}%
\usepackage[font=small,labelfont=bf]{caption}%
%

            \hypersetup{
                colorlinks=true,
                linkcolor=blue,
                filecolor=magenta,
                urlcolor=blue,
                breaklinks=true,
                bookmarks=true,
                pdfborder={0 0 0}
            }
            \urlstyle{same}
            \Urlmuskip=0mu plus 1mu  % URL breaking in nicer spots

            % Math configuration
            \everymath{\displaystyle}
            \setlength{\jot}{10pt}  % Increase spacing between equations

            \pagestyle{fancy}
            \fancyhf{}
            \rhead{Generated on \today}
            \cfoot{\thepage}

            % Better Unicode support
            \usepackage[utf8]{inputenc}
            \usepackage{newunicodechar}
            \setlength{\itemsep}{1em} % Adds spacing between items

            % Configure itemize settings globally
            \setlist[itemize]{nosep}
            \setlist[itemize]{leftmargin=*}

            % Remove section numbering
            \renewcommand{\thesection}{}
            \renewcommand{\thesubsection}{}

            % Adjust caption spacing around figures
            \setlength{\abovecaptionskip}{0pt}
            \setlength{\belowcaptionskip}{5pt}
        %
\lhead{Lecture: 12L2{-}distance{-}clusters}%
\title{12L2{-}distance{-}clusters}%
\author{Generated by Scribe.AI}%
\date{\today}%
%
\begin{document}%
\normalsize%
\maketitle%
\section{Slide 1}%
\label{sec:Slide1}%
\subsection{Content}%

\textbf{AI Basics} \\
Purdue CS24300 \\
Tony Bergstrom \\
Slides adapted from Jen Neville, CS 24200
%
\subsection{Description}%
This slide is a title slide for a lecture on Artificial Intelligence Basics.  It displays the course title, "AI Basics," the course number, "Purdue CS24300," the instructor's name, "Tony Bergstrom," and an attribution line indicating that the slides are adapted from another source, "Slides adapted from Jen Neville, CS 24200."  The text is presented in a simple, clear format on a dark gray background.  The context of the course is evident from the title and course number, indicating that this is the introductory slide to a lecture series on fundamental concepts in Artificial Intelligence.%
\subsection{Figures}%
\begin{figure}[H]%
\centering%
\end{figure}

%
\newpage%
\section{Slide 2}%
\label{sec:Slide2}%
\subsection{Content}%

\textbf{Introduction to Clustering} \\
Data Science in Business, Chapter 6
%
\subsection{Description}%
This slide displays the title "Introduction to Clustering" in a large, bold font.  Below the title, in a smaller font, it states "Data Science in Business, Chapter 6."  The background is a neutral gray color.  The context of the slide is an introductory lecture on clustering techniques within a data science course.  The phrase "Data Science in Business" indicates that the course is focused on the application of data science principles in a business setting. The "Chapter 6" designation suggests that this is part of a larger series of lectures on data science topics.%
\subsection{Figures}%
\begin{figure}[H]%
\centering%
\end{figure}

%
\newpage%
\section{Slide 3}%
\label{sec:Slide3}%
\subsection{Content}%

\textbf{Subfields of Machine learning} \\
\textbf{Supervised Learning} \\
Given paired inputs/outputs (X, y) \\
Goal is to learn $f$ from training data, $f(X) = y$ \\
Evaluated it on "new" (test) data. \\
Examples \\
\begin{itemize}
\item Classification: y is discrete (y are class labels, X refers to attributes)
\item Regression: y is continuous
\end{itemize}
\textbf{Unsupervised learning} \\
Given examples of inputs X \\
Learn a function $f$ to simplify the data, $f(X) = y$ \\
Examples: \\
\begin{itemize}
\item Clustering: y is discrete (y are cluster memberships)
\item Matrix factorization: y are continuous (y are embeddings)
\end{itemize}
%
\subsection{Description}%
This slide presents a breakdown of machine learning subfields, specifically focusing on supervised and unsupervised learning.  The slide is titled "Subfields of Machine Learning."  Under "Supervised Learning," it details that paired inputs and outputs (X, y) are given, and the goal is to learn a function $f$ from training data such that $f(X) = y$.  The learned function is then evaluated on new, unseen data.  Examples of supervised learning are listed: classification (where the output y is discrete, representing class labels) and regression (where the output y is continuous).  Under "Unsupervised Learning," the slide describes that only input examples X are given, and the goal is to learn a function $f$ that simplifies the data, again with $f(X) = y$.  Examples of unsupervised learning are clustering (where the output y is discrete, representing cluster memberships) and matrix factorization (where the output y is continuous, representing embeddings). The context of the course (Artificial Intelligence Basics) indicates that this slide is introducing fundamental concepts in machine learning, specifically the different types of learning paradigms.%
\subsection{Figures}%
\begin{figure}[H]%
\centering%
\end{figure}

%
\newpage%
\section{Slide 4}%
\label{sec:Slide4}%
\subsection{Content}%

\textbf{What is Clustering?} \\
Clustering: the process of grouping a set of objects into classes of similar objects \\
\begin{itemize}
\item Objects within a cluster should be similar
\item Objects from different clusters should be dissimilar
\end{itemize}
Clustering is the most common form of unsupervised learning \\
\begin{itemize}
\item Learning from unlabeled raw data
\item A frequent and important task that finds many applications in data science
\end{itemize}
%
\subsection{Description}%
This slide introduces the concept of clustering.  It defines clustering as the process of grouping a set of objects into classes of similar objects.  Key characteristics of a good clustering are that objects within a cluster should be similar to each other, and objects from different clusters should be dissimilar.  The slide further emphasizes that clustering is a common form of unsupervised learning, meaning it learns from unlabeled data, and that it is a frequently used technique in data science applications.%
\subsection{Figures}%
\begin{figure}[H]%
\centering%
\end{figure}

%
\newpage%
\section{Slide 5}%
\label{sec:Slide5}%
\subsection{Content}%

\textbf{Examples}
%
\subsection{Description}%
This slide, titled "Examples," is likely a placeholder or a slide that will contain examples related to the previous topics discussed in the lecture.  The gray background and the single word "Examples" suggest that this slide is empty and will be filled with content later, possibly diagrams, graphs, or descriptions of specific examples relevant to the course, "Artificial Intelligence Basics."  The context of the course is important to understand the type of examples that will be presented.  For example, if the previous slides discussed clustering algorithms, the examples might show different datasets and how the algorithms would group them.%
\subsection{Figures}%
\begin{figure}[H]%
\centering%
\end{figure}

%
\newpage%
\section{Slide 6}%
\label{sec:Slide6}%
\subsection{Content}%

\textbf{Example: Documents}
%
\subsection{Description}%
This slide, titled "Example: Documents," displays a collection of news articles or press releases.  The content appears to be a mix of articles from various news sources, likely from different countries (U.S. and World).  The articles cover a range of topics, including political events (e.g., elections, political statements), social issues (e.g., anti-Biden chants), and world events (e.g., a man dressed as the Joker injuring people on a train).  The slide likely serves as an example of the types of documents that might be used as data in an AI project, such as sentiment analysis, topic modeling, or news aggregation.  The context of the course, "Artificial Intelligence Basics," suggests that this slide is meant to illustrate how real-world data can be used as input for AI models.  The presence of images and links to "View Full Coverage" indicates that the articles are likely from online news sources.  The slide does not contain any specific instructions or analysis, but rather presents a collection of examples for potential use in AI projects.%
\subsection{Figures}%
\begin{figure}[H]%
\centering%
\begin{subfigure}[b]{0.15\textwidth}%
\centering%
\includegraphics[width=\textwidth]{/Users/ashoksaravanan/Coding/ScribeLec/server/output/CS 243/lectures/12L2-distance-clusters/figures/6.1.png}%
\caption{}%
\end{subfigure}%
\hspace{0.05\textwidth}%
\begin{subfigure}[b]{0.15\textwidth}%
\centering%
\includegraphics[width=\textwidth]{/Users/ashoksaravanan/Coding/ScribeLec/server/output/CS 243/lectures/12L2-distance-clusters/figures/6.2.png}%
\caption{}%
\end{subfigure}%
\hspace{0.05\textwidth}%
\begin{subfigure}[b]{0.15\textwidth}%
\centering%

\includegraphics[width=\textwidth]{/Users/ashoksaravanan/Coding/ScribeLec/server/output/CS 243/lectures/12L2-distance-clusters/figures/6.3.png}%
\caption{}%
\end{subfigure}%
\end{figure}

%
\newpage%
\section{Slide 7}%
\label{sec:Slide7}%
\subsection{Content}%

\textbf{Example: Images}
%
\subsection{Description}%
This slide, titled "Example: Images," presents an example of image segmentation.  It shows a street scene image (likely a photograph) at the top, and a segmented version of the same image below it.  The segmented image is a color-coded representation of the scene, where different colors correspond to different object categories.  Labels such as "STATIONARY VEGETATION," "MOVING CAR," "STATIONARY ROAD," "STATIONARY BUILDING," and "STATIONARY PAVEMENT" are shown next to the corresponding colored regions in the segmented image.  The labels indicate that the image segmentation algorithm has identified and categorized different objects in the scene.  The context of the course, "Artificial Intelligence Basics," suggests that this slide is illustrating a computer vision task, specifically image segmentation, which is a crucial component of many AI applications, such as self-driving cars, medical image analysis, and robotics.  The segmented image shows the algorithm's attempt to identify and classify different objects in the scene, highlighting the potential of AI to understand and interpret visual data.%
\subsection{Figures}%
\begin{figure}[H]%
\centering%
\begin{subfigure}[b]{0.15\textwidth}%
\centering%
\includegraphics[width=\textwidth]{/Users/ashoksaravanan/Coding/ScribeLec/server/output/CS 243/lectures/12L2-distance-clusters/figures/7.1.png}%
\caption{}%
\end{subfigure}%
\end{figure}

%
\newpage%
\section{Slide 8}%
\label{sec:Slide8}%
\subsection{Content}%

\textbf{Example: Genes}
%
\subsection{Description}%
This slide, titled "Example: Genes," displays a complex network graph.  The nodes in the graph represent genes or gene sets, and the connections between them indicate relationships or co-expression patterns.  Different colors are used to visually distinguish various cell types or tissues, such as "Astrocytoma," "Smooth muscle," "Chorion, placenta, choriocarcinoma cell lines," "Epithelial cells and tissues," "Mesenchymal cells, fibroblasts," "Cardia," "Skeletal muscle," "Lymphoma cell lines, reticulocytes," "Leukemia/lymphoma cell lines," "Lymphoid," "Retina, melanocytes, melanoma," and others.  The graph visually represents a promoter-level mammalian expression atlas, showing how different genes are expressed in various cell types.  The context of the course, "Artificial Intelligence Basics," suggests that this slide is illustrating a biological dataset that could be used as input for AI models to analyze gene expression patterns, predict disease states, or discover relationships between genes and cell types.  The slide also includes labels and annotations for specific genes or pathways, indicating the biological significance of the connections.  The citation "A promoter-level mammalian expression atlas. Nature, 26 March 2014" at the bottom of the slide indicates the source of the data.  The visual complexity of the graph highlights the vast amount of data and the intricate relationships within gene expression.%
\subsection{Figures}%
\begin{figure}[H]%
\centering%
\begin{subfigure}[b]{0.15\textwidth}%
\centering%
\includegraphics[width=\textwidth]{/Users/ashoksaravanan/Coding/ScribeLec/server/output/CS 243/lectures/12L2-distance-clusters/figures/8.1.png}%
\caption{}%
\end{subfigure}%
\end{figure}

%
\newpage%
\section{Slide 9}%
\label{sec:Slide9}%
\subsection{Content}%

\textbf{Example: Movies}
%
\subsection{Description}%
This slide, titled "Example: Movies," presents a visualization of 17,700 movies from the Netflix competition.  It's a scatter plot-like representation, where each small dot represents a movie.  Different colors and clusters of dots visually group movies based on various attributes or characteristics.  Labels such as "Exercise," "Bonus Material," "Family," "Super Size," "Critically Acclaimed," "Sci-Fi," "The Passion of the Christ," "The Ring," "Battlestar Galactica," "A Beautiful Mind," "When We Were Kings," "Sin City," and "Brave Heart" are placed near clusters of movies, indicating the categories or themes associated with those groups.  The slide also shows the names of individual movies, such as "G.I. Jane," "Armageddon," "Independence Day," "Dark Water," "Dead Birds," "Equilibrium," "The Jacket," "Leperchaun 4 in Space," and "Crocodile Dundee."  The context of the course, "Artificial Intelligence Basics," suggests that this slide is illustrating a dataset of movies and their associated features, potentially for use in AI tasks such as movie recommendation systems, genre classification, or clustering movies based on similar characteristics.  The visualization highlights the vast amount of data and the potential for AI to analyze and categorize such large datasets.  The date "03/25/2007" and email address "Todd.Holloway@gmail.com" suggest the source and creation date of the data or visualization.  The presence of movie titles and categories indicates that the data is likely from a movie database or a similar source.%
\subsection{Figures}%
\begin{figure}[H]%
\centering%
\begin{subfigure}[b]{0.15\textwidth}%
\centering%
\includegraphics[width=\textwidth]{/Users/ashoksaravanan/Coding/ScribeLec/server/output/CS 243/lectures/12L2-distance-clusters/figures/9.1.png}%
\caption{}%
\end{subfigure}%
\end{figure}

%
\newpage%
\section{Slide 10}%
\label{sec:Slide10}%
\subsection{Content}%

\textbf{Phylogenetic Tree of Dog Breeds}
%
\subsection{Description}%
This slide displays a circular phylogenetic tree, visualizing the evolutionary relationships between various dog breeds.  The tree structure shows the branching patterns, with nodes representing common ancestors and branches connecting different breeds.  Each breed is represented by a colored dot, with different colors likely indicating different clusters or groups of breeds.  The colors are not explicitly labeled, but the legend in the bottom right corner shows that different sizes of dots correspond to different ranges of similarity scores (e.g., gold dots = 90-100, dark gray dots = 50-69).  The names of numerous dog breeds are labeled along the branches, including breeds like "West Highland White Terrier," "German Shepherd," "Alaskan Malamute," "Siberian Husky," "Boxer," "Bull Terrier," "Great Dane," and many others.  The circular layout allows for a comprehensive view of the relationships between a large number of dog breeds.  The context of the course, "Artificial Intelligence Basics," suggests that this slide is illustrating a dataset of dog breeds and their genetic or phenotypic similarities, potentially for use in AI tasks such as breed classification, prediction of traits, or analysis of genetic diversity.  The visualization highlights the vast amount of data and the potential for AI to analyze and categorize such large datasets.  The legend indicates that the data is likely based on similarity scores, which could be derived from genetic or phenotypic data.%
\subsection{Figures}%
\begin{figure}[H]%
\centering%
\begin{subfigure}[b]{0.15\textwidth}%
\centering%
\includegraphics[width=\textwidth]{/Users/ashoksaravanan/Coding/ScribeLec/server/output/CS 243/lectures/12L2-distance-clusters/figures/10.1.png}%
\caption{}%
\end{subfigure}%
\hspace{0.05\textwidth}%
\begin{subfigure}[b]{0.15\textwidth}%
\centering%
\includegraphics[width=\textwidth]{/Users/ashoksaravanan/Coding/ScribeLec/server/output/CS 243/lectures/12L2-distance-clusters/figures/10.2.png}%
\caption{}%
\end{subfigure}%
\end{figure}

%
\newpage%
\section{Slide 11}%
\label{sec:Slide11}%
\subsection{Content}%

\textbf{Personas}
%
\subsection{Description}%
This slide, titled "Personas," displays a row of twelve simple, cartoon-like portraits of diverse individuals.  The individuals represent different genders, ages, and ethnicities.  The style is a basic outline drawing, with colors used to differentiate the individuals.  The text below the images states, "To provide a common understanding across all teams."  The context of the course, "Artificial Intelligence Basics," suggests that these personas are likely part of a design thinking or user-centered design process.  The personas are intended to represent different user types or target audiences for a project or product.  The goal is to create a shared understanding of these user types across different teams working on the project.  The lack of specific details about the personas (e.g., names, occupations, motivations) indicates that this is a high-level overview, likely used as a starting point for further development and refinement of user profiles.%
\subsection{Figures}%
\begin{figure}[H]%
\centering%
\begin{subfigure}[b]{0.15\textwidth}%
\centering%
\includegraphics[width=\textwidth]{/Users/ashoksaravanan/Coding/ScribeLec/server/output/CS 243/lectures/12L2-distance-clusters/figures/11.1.png}%
\caption{}%
\end{subfigure}%
\hspace{0.05\textwidth}%
\begin{subfigure}[b]{0.15\textwidth}%
\centering%

\includegraphics[width=\textwidth]{/Users/ashoksaravanan/Coding/ScribeLec/server/output/CS 243/lectures/12L2-distance-clusters/figures/11.2.png}%
\caption{}%
\end{subfigure}%
\end{figure}

%
\newpage%
\section{Slide 12}%
\label{sec:Slide12}%
\subsection{Content}%
%
\subsection{Description}%
This slide, titled "Gathering Data," presents a summary of data collection efforts.  It lists the number of interviews conducted with different groups (Instructors, Instructional Designers, Administrators, Students, and Coursona).  Each group is associated with a count and a breakdown of the number of participants from the United States (U.S.).  The total number of interviews is 90, each lasting one hour, and 81% of the participants were from the U.S.  The data likely represents a study or research project focused on understanding the perspectives and experiences of these groups.  The context of the course, "Artificial Intelligence Basics," suggests that this data might be used to inform or train AI models, potentially for tasks like sentiment analysis, user profiling, or understanding user needs in educational settings.  The table format, with the breakdown of U.S. participants, implies that the data is being analyzed for potential cultural or regional differences in perspectives.%
\subsection{Figures}%
\begin{figure}[H]%
\centering%
\begin{subfigure}[b]{0.15\textwidth}%
\centering%
\includegraphics[width=\textwidth]{/Users/ashoksaravanan/Coding/ScribeLec/server/output/CS 243/lectures/12L2-distance-clusters/figures/12.1.png}%
\caption{}%
\end{subfigure}%
\hspace{0.05\textwidth}%
\begin{subfigure}[b]{0.15\textwidth}%
\centering%
\includegraphics[width=\textwidth]{/Users/ashoksaravanan/Coding/ScribeLec/server/output/CS 243/lectures/12L2-distance-clusters/figures/12.2.png}%
\caption{}%
\end{subfigure}%
\end{figure}

%
\newpage%
\section{Slide 13}%
\label{sec:Slide13}%
\subsection{Content}%

\textbf{Grouping the data}

Association Analysis \\
Known Separation \\
Hierarchical Cluster Analysis
%
\subsection{Description}%
This slide, titled "Grouping the data," lists three methods for analyzing and grouping data: Association Analysis, Known Separation, and Hierarchical Cluster Analysis.  The context of the course, "Artificial Intelligence Basics," suggests that these are different techniques used in data mining and machine learning to discover patterns and relationships within datasets.  The slide likely precedes a more detailed explanation of each method, possibly including examples, algorithms, or applications.  The image in the right-hand corner shows a group of people working together, likely brainstorming or collaborating on a project, which could be related to the data analysis process.  The image of the data grouping methods below the image of the people is a visual representation of hierarchical clustering, showing how different data points are grouped together in a tree-like structure.  The colors represent different clusters, and the heights of the bars represent the similarity or distance between the data points.  The slide is likely part of a section on data preprocessing or feature engineering, where the goal is to organize and prepare data for use in AI algorithms.%
\subsection{Figures}%
\begin{figure}[H]%
\centering%
\begin{subfigure}[b]{0.15\textwidth}%
\centering%
\includegraphics[width=\textwidth]{/Users/ashoksaravanan/Coding/ScribeLec/server/output/CS 243/lectures/12L2-distance-clusters/figures/13.1.png}%
\caption{}%
\end{subfigure}%
\end{figure}

%
\newpage%
\section{Slide 14}%
\label{sec:Slide14}%
\subsection{Content}%
%
\subsection{Description}%
This slide, titled "Building a Narrative," presents a persona profile for Patti Caresalot, a full-time English instructor at Lake Valley College.  The persona is depicted in a card-like format, with sections for "Higher Ed.", "Description", "Challenges", and "Personal".  The "Higher Ed." section describes Patti's teaching priorities, focusing on student engagement and actively seeking feedback.  The "Description" section highlights Patti's belief that learning is most effective when students actively apply knowledge and engage in discussions.  The "Challenges" section details the instructor's current workload, teaching four courses, two online and two in-person, with varying class sizes (30 and 90 students respectively).  The "Personal" section likely contains information about Patti's personal characteristics or preferences, which are not visible in the provided image.  The "Coursona Matches" section likely lists matching criteria or characteristics based on a platform or system called Coursona, used for identifying and categorizing user profiles.  The overall context suggests this is part of a user-centered design or instructional design process, focusing on understanding different instructor profiles and their needs.%
\subsection{Figures}%
\begin{figure}[H]%
\centering%
\begin{subfigure}[b]{0.15\textwidth}%
\centering%
\includegraphics[width=\textwidth]{/Users/ashoksaravanan/Coding/ScribeLec/server/output/CS 243/lectures/12L2-distance-clusters/figures/14.1.png}%
\caption{}%
\end{subfigure}%
\end{figure}

%
\newpage%
\section{Slide 15}%
\label{sec:Slide15}%
\subsection{Content}%

\textbf{Clustering terminology}

\textbf{Hierarchical clustering}
\begin{itemize}
    \item clusters form a hierarchy
    \item computed bottom-up or top-down fashion
\end{itemize}

\textbf{Flat clustering}
\begin{itemize}
    \item no inter-cluster structure
\end{itemize}

\textbf{Hard clustering}
\begin{itemize}
    \item items are assigned to a single unique cluster
\end{itemize}

\textbf{Soft clustering}
\begin{itemize}
    \item items are assigned to one or more clusters
    \item membership often expressed as a real-valued function
\end{itemize}

Kmeans is flat, hard clustering approach
%
\subsection{Description}%
This slide, titled "Clustering terminology," defines different types of clustering methods.  It presents a table-like format, listing the characteristics of hierarchical, flat, hard, and soft clustering.  The context of the course, "Artificial Intelligence Basics," suggests that this slide is part of a module on unsupervised machine learning techniques.  Hierarchical clustering is described as creating a hierarchy of clusters, computed either bottom-up or top-down.  Flat clustering, in contrast, has no inter-cluster structure, meaning clusters are not organized in a hierarchical manner.  Hard clustering assigns each data point to a single cluster, while soft clustering allows data points to belong to multiple clusters with varying degrees of membership, often represented by a real-valued function.  The final line, "Kmeans is flat, hard clustering approach," highlights a specific clustering algorithm, K-means, as an example of a flat and hard clustering method.  This slide is likely intended to provide a foundational understanding of different clustering approaches before delving into specific algorithms or applications.%
\subsection{Figures}%
\begin{figure}[H]%
\centering%
\end{figure}

%
\newpage%
\section{Slide 16}%
\label{sec:Slide16}%
\subsection{Content}%

\textbf{Similarity/Distance}
%
\subsection{Description}%
This slide, likely part of a presentation on Artificial Intelligence Basics, displays a single heading: "Similarity/Distance."  It indicates that the subsequent content will focus on methods for measuring similarity or distance between data points.  The context suggests this is a topic within data preprocessing, feature engineering, or clustering algorithms, where determining the similarity or dissimilarity of data is crucial for effective analysis.%
\subsection{Figures}%
\begin{figure}[H]%
\centering%
\end{figure}

%
\newpage%
\section{Slide 17}%
\label{sec:Slide17}%
\subsection{Content}%

\textbf{Distance measures}

If objects have the same fixed set of numeric attributes, then the
data objects can be thought of as points in a multi-dimensional
space, where each dimension represents a distinct attribute

Euclidean distance is the most commonly used
distance measure for continuous data:

$d(x, y) = \sqrt{\sum_{i=1}^{p} (x_i - y_i)^2}$

Here $p$ is the number of attributes/dimensions

Proximity: general term to indicate similarity
and dissimilarity

Distance: measure of dissimilarity

Heatmap of distances between 100 different elements
%
\subsection{Description}%
This slide, titled "Distance measures," discusses methods for quantifying the similarity or dissimilarity between data points in a multi-dimensional space.  It defines Euclidean distance as the most common measure for continuous data.  The formula for Euclidean distance is presented: $d(x, y) = \sqrt{\sum_{i=1}^{p} (x_i - y_i)^2}$, where $x$ and $y$ represent data points, $x_i$ and $y_i$ are the values of the $i$-th attribute for points $x$ and $y$ respectively, and $p$ is the total number of attributes.  The slide also introduces the terms "Proximity" and "Distance," which are general terms used to describe the degree of similarity or dissimilarity between data points.  A heatmap visualization of distances between 100 different elements is shown, illustrating the application of these distance measures.  The context of the course, "Artificial Intelligence Basics," suggests that this slide is part of a module on data preprocessing, feature engineering, or clustering algorithms, where understanding distance metrics is crucial for comparing and grouping data points.%
\subsection{Figures}%
\begin{figure}[H]%
\centering%
\begin{subfigure}[b]{0.15\textwidth}%
\centering%
\includegraphics[width=\textwidth]{/Users/ashoksaravanan/Coding/ScribeLec/server/output/CS 243/lectures/12L2-distance-clusters/figures/17.1.png}%
\caption{}%
\end{subfigure}%
\end{figure}

%
\newpage%
\section{Slide 18}%
\label{sec:Slide18}%
\subsection{Content}%
%
\subsection{Description}%
This slide demonstrates an example of using the matplotlib library in Python, likely within a course on Artificial Intelligence Basics.  The code snippet calculates a distance matrix using the `distance_matrix` function from the `scipy.spatial` library.  The `xyLocations` variable generates 25 data points, each with two coordinates: an integer between 0 and 9, and a Gaussian-distributed value centered around 5 with a standard deviation of 3.  These points are then stored in a Pandas DataFrame called `xyLocs`.  The `distance_matrix` function computes the Euclidean distances between all pairs of points in `xyLocs`.  The `print(distances[:5, :5])` command displays the top-left 5x5 portion of the distance matrix, showing the distances between the first five data points.  The `plt.imshow(distances)` command then displays a heatmap visualization of the entire distance matrix.  The heatmap shows the distances between all 25 data points, with darker colors representing smaller distances and lighter colors representing larger distances.  The context suggests that this slide is part of a section on data visualization and distance metrics, which are crucial for various AI tasks, such as clustering and pattern recognition.%
\subsection{Figures}%
\begin{figure}[H]%
\centering%
\begin{subfigure}[b]{0.15\textwidth}%
\centering%
\includegraphics[width=\textwidth]{/Users/ashoksaravanan/Coding/ScribeLec/server/output/CS 243/lectures/12L2-distance-clusters/figures/18.1.png}%
\caption{}%
\end{subfigure}%
\end{figure}

%
\newpage%
\section{Slide 19}%
\label{sec:Slide19}%
\subsection{Content}%

\textbf{Distance between discrete data}

\textbf{Hamming distance}
\begin{itemize}
    \item Used for bit/binary vectors: counts number of bits that differ
    \item Can be used more generally for discrete data, to count matches
\end{itemize}

\textbf{Matching coefficient}
\begin{itemize}
    \item Hamming distance normalized by number of attributes/bits
\end{itemize}

A  1 0 1 1 1 0 0 1 0 0 1
B  1 0 0 1 0 0 0 0 1 1 1

Hamming Distance: 3
Matching coefficient: 7/10
%
\subsection{Description}%
This slide discusses distance measures for discrete data, specifically focusing on Hamming distance and matching coefficient.  It explains that Hamming distance counts the number of differing bits (or attributes) between two bit/binary vectors.  The example shows two vectors, A and B, each with 11 bits.  The Hamming distance between A and B is calculated as 3, because there are 3 positions where the bits differ.  The matching coefficient is calculated as 7/11, which represents the proportion of matching bits (7 matches out of 10 total bits).  The slide highlights that Hamming distance can be applied more broadly to discrete data, not just binary vectors, to count the number of mismatches.  The matching coefficient normalizes the Hamming distance by dividing it by the total number of attributes/bits.  The context of the course, "Artificial Intelligence Basics," suggests that this slide is part of a module on data preprocessing, feature engineering, or similarity measures, where understanding different distance metrics for various data types is crucial.%
\subsection{Figures}%
\begin{figure}[H]%
\centering%
\begin{subfigure}[b]{0.15\textwidth}%
\centering%
\includegraphics[width=\textwidth]{/Users/ashoksaravanan/Coding/ScribeLec/server/output/CS 243/lectures/12L2-distance-clusters/figures/19.1.png}%
\caption{}%
\end{subfigure}%
\end{figure}

%
\newpage%
\section{Slide 20}%
\label{sec:Slide20}%
\subsection{Content}%

\textbf{Clustering}
%
\subsection{Description}%
This slide, likely part of a lecture on Artificial Intelligence Basics, simply displays the title "Clustering."  It indicates that the following slides will cover topics related to clustering algorithms, a common technique in machine learning and data mining used to group similar data points together.%
\subsection{Figures}%
\begin{figure}[H]%
\centering%
\end{figure}

%
\newpage%
\section{Slide 21}%
\label{sec:Slide21}%
\subsection{Content}%

\textbf{K-means clustering}

One of the most commonly-used clustering algorithms, because it is easy to implement
and quick to run

The goal is to partition the objects into $K$ disjoint subsets

Assumes the objects (instances) to be clustered are $p$-dimensional vectors

Uses a distance measure between the objects (typically Euclidean distance)
%
\subsection{Description}%
This slide introduces the K-means clustering algorithm, a common technique in unsupervised machine learning.  It states that K-means is one of the most frequently used clustering algorithms due to its ease of implementation and speed.  The goal of K-means is to divide the input data points into $K$ distinct, non-overlapping groups.  The algorithm assumes that the data points to be clustered are represented as vectors in a $p$-dimensional space.  Finally, it specifies that a distance measure, typically the Euclidean distance, is used to determine the similarity between data points for grouping them into clusters.  The context of the course, "Artificial Intelligence Basics," suggests that this slide is part of a module on clustering algorithms and their applications in various AI tasks.%
\subsection{Figures}%
\begin{figure}[H]%
\centering%
\end{figure}

%
\newpage%
\section{Slide 22}%
\label{sec:Slide22}%
\subsection{Content}%

\textbf{The K-means Algorithm}

\textbf{Input:} Dataset $D$, number of clusters $k$

\textbf{Output:} Set of cluster representatives $C$, cluster membership vector $m$

\textbf{\# Initialize cluster representatives $C$}

Randomly assign $k$ data points from $D$ into each cluster

Initialize $C$ such that $c_i$ is mean of points in $i^{th}$ cluster

\textbf{while} not converged

\textbf{\# Data Assignment}

Reassign points in $D$ to closest cluster mean

Update $m$ such that $m_i$ is cluster ID of the $i^{th}$ point in $D$

\textbf{\# Relocation of means}

Update $C$ such that $c_i$ is mean of points in $i^{th}$ cluster
%
\subsection{Description}%
This slide presents the K-means clustering algorithm.  It outlines the steps involved in the algorithm, clearly defining the input (a dataset $D$ and the number of clusters $k$) and the output (a set of cluster representatives $C$ and a cluster membership vector $m$).  The algorithm begins by initializing the cluster representatives $C$ by randomly selecting $k$ data points from the dataset $D$ to be the initial cluster centers.  The algorithm then iteratively performs two steps:  Data Assignment and Relocation of means.  In the Data Assignment step, each data point in $D$ is reassigned to the closest cluster mean.  The cluster membership vector $m$ is updated to reflect these assignments.  In the Relocation of means step, the cluster representatives $C$ are updated by recalculating the mean of the data points assigned to each cluster.  The algorithm continues these steps until it converges, meaning that the cluster assignments no longer change significantly between iterations.  The context of the course, "Artificial Intelligence Basics," suggests that this slide is part of a module on clustering algorithms and their applications in various AI tasks.%
\subsection{Figures}%
\begin{figure}[H]%
\centering%
\begin{subfigure}[b]{0.15\textwidth}%
\centering%

\includegraphics[width=\textwidth]{/Users/ashoksaravanan/Coding/ScribeLec/server/output/CS 243/lectures/12L2-distance-clusters/figures/22.1.png}%
\caption{}%
\end{subfigure}%
\end{figure}

%
\newpage%
\section{Slide 23}%
\label{sec:Slide23}%
\subsection{Content}%
%
\subsection{Description}%
This slide, likely part of a lecture on clustering algorithms within the context of Artificial Intelligence Basics, displays a scatter plot.  The plot shows a set of data points, represented as small black dots, scattered across a two-dimensional space.  The title "Example: Acquire Data" suggests that this is an example of data acquisition, illustrating the raw data that a clustering algorithm might receive as input.  The data points appear to be randomly distributed, without any discernible clusters or patterns.  The context of the course indicates that this slide is meant to show the type of data that clustering algorithms are designed to analyze and group into clusters.%
\subsection{Figures}%
\begin{figure}[H]%
\centering%
\end{figure}

%
\newpage%
\section{Slide 24}%
\label{sec:Slide24}%
\subsection{Content}%

\textbf{Example: Assign Randomly}
%
\subsection{Description}%
This slide displays a scatter plot, likely part of a lecture on the K-means clustering algorithm within the context of Artificial Intelligence Basics.  The plot shows a set of data points, represented as colored circles (red, gray, and light blue), randomly scattered across a two-dimensional space.  The title "Example: Assign Randomly" indicates that this is an example of the initial random assignment of data points to clusters in the K-means algorithm.  Each color represents a different cluster, and the points are not grouped into any discernible clusters.  The context of the course suggests that this slide is meant to illustrate the initial state of the data before the K-means algorithm begins iteratively refining the cluster assignments.%
\subsection{Figures}%
\begin{figure}[H]%
\centering%
\end{figure}

%
\newpage%
\section{Slide 25}%
\label{sec:Slide25}%
\subsection{Content}%

\textbf{Example: Compute Centroids}
%
\subsection{Description}%
This slide, likely part of a lecture on the K-means clustering algorithm within the context of Artificial Intelligence Basics, displays a scatter plot.  The plot shows a set of data points, represented as small filled circles in various colors (red, gray, and light blue).  The title "Example: Compute Centroids" indicates that this is an example of the centroid calculation step in the K-means algorithm.  The data points are scattered across a two-dimensional space, and several points are marked with larger, differently colored stars.  These stars represent the computed centroids for the clusters.  The context of the course suggests that this slide is meant to illustrate the result of the K-means algorithm after the initial random assignment and relocation of cluster means.  The different colored points represent different clusters, and the corresponding centroids are marked with stars of the same color.  The slide visually demonstrates the concept of finding the mean of the data points assigned to each cluster to determine the cluster centroids.%
\subsection{Figures}%
\begin{figure}[H]%
\centering%
\end{figure}

%
\newpage%
\section{Slide 26}%
\label{sec:Slide26}%
\subsection{Content}%

\textbf{Example: Reassign Clusters}
%
\subsection{Description}%
This slide, likely part of a lecture on the K-means clustering algorithm within the context of Artificial Intelligence Basics, displays a scatter plot.  The plot shows a set of data points, represented as small filled circles in three different colors (red, gray, and light blue).  The title "Example: Reassign Clusters" indicates that this is an example of the iterative reassignment of data points to clusters in the K-means algorithm.  Each color represents a different cluster.  The data points are scattered across a two-dimensional space.  Importantly, the points are not randomly assigned as in a previous slide, but rather, the points are grouped into clusters.  Several points are marked with larger, differently colored stars.  These stars represent the centroids for the clusters.  The slide visually demonstrates the intermediate step in the K-means algorithm where data points are reassigned to the nearest cluster centroid.  The context of the course suggests that this slide is meant to illustrate the algorithm's iterative process of refining cluster assignments until convergence.  The different colored points represent different clusters, and the corresponding centroids are marked with stars of the same color.  The slide shows the result of the K-means algorithm after multiple iterations, where the data points are more clearly grouped into clusters.%
\subsection{Figures}%
\begin{figure}[H]%
\centering%
\end{figure}

%
\newpage%
\section{Slide 27}%
\label{sec:Slide27}%
\subsection{Content}%

\textbf{Example: Re-compute Centroids}
%
\subsection{Description}%
This slide, part of a lecture on K-means clustering within an Artificial Intelligence Basics course, displays a scatter plot.  The plot shows data points clustered into three groups, represented by small filled circles in three different colors:  red, light blue, and gray.  The title "Example: Re-compute Centroids" indicates that this is a step in the K-means algorithm.  The key visual element is the presence of three larger, differently colored stars. These stars represent the centroids for the three clusters.  The red star is the centroid for the red cluster, the light blue star is the centroid for the light blue cluster, and the gray star is the centroid for the gray cluster.  The data points are scattered across a two-dimensional space, and the distribution of points around each centroid suggests that the algorithm has already performed one or more iterations of cluster assignment and centroid recalculation.  The context of the course suggests that this slide is meant to illustrate the algorithm's iterative process of refining cluster assignments and centroids until convergence.  The different colored points represent different clusters, and the corresponding centroids are marked with stars of the same color.  The slide shows the result of the K-means algorithm after an intermediate step, where the centroids have been recalculated based on the current cluster assignments.%
\subsection{Figures}%
\begin{figure}[H]%
\centering%
\end{figure}

%
\newpage%
\section{Slide 28}%
\label{sec:Slide28}%
\subsection{Content}%

\textbf{Example: Reassign Clusters}
%
\subsection{Description}%
This slide, part of a lecture on K-means clustering in an Artificial Intelligence Basics course, displays a scatter plot.  The plot shows data points clustered into three groups, represented by small filled circles in three different colors: red, light blue, and gray.  The title "Example: Reassign Clusters" indicates that this is a step in the K-means algorithm.  The data points are scattered across a two-dimensional space.  The slide visually demonstrates the intermediate step in the K-means algorithm where data points are reassigned to the nearest cluster centroid.  The different colored points represent different clusters.  Crucially, the slide shows the result *after* the initial random assignment of points to clusters, and *before* the centroids are recomputed.  The positions of the points relative to the centroids (which are not explicitly labeled but are implied by the color groupings) indicate that the algorithm has already performed one or more iterations of cluster assignment.  The context of the course suggests that this slide is meant to illustrate the iterative process of refining cluster assignments until convergence.  The different colored points represent different clusters.  The slide does not show the centroids themselves, but the color-coding of the data points implies that the algorithm has already grouped the data into clusters.%
\subsection{Figures}%
\begin{figure}[H]%
\centering%
\end{figure}

%
\newpage%
\section{Slide 29}%
\label{sec:Slide29}%
\subsection{Content}%

\textbf{Example: Re-compute Centroids}
%
\subsection{Description}%
This slide, part of a lecture on K-means clustering within an Artificial Intelligence Basics course, displays a scatter plot.  The plot shows data points clustered into three groups, represented by small filled circles in three different colors: red, light blue, and gray.  The title "Example: Re-compute Centroids" indicates that this is a step in the K-means algorithm.  The key visual element is the presence of three larger, differently colored stars. These stars represent the centroids for the three clusters.  The red star is the centroid for the red cluster, the light blue star is the centroid for the light blue cluster, and the gray star is the centroid for the gray cluster.  The data points are scattered across a two-dimensional space, and the distribution of points around each centroid suggests that the algorithm has already performed one or more iterations of cluster assignment.  The slide is illustrating the K-means algorithm's iterative process of refining cluster assignments and centroids until convergence.  The different colored points represent different clusters, and the corresponding centroids are marked with stars of the same color.  The slide shows the result of the K-means algorithm after an intermediate step, where the centroids have been recalculated based on the current cluster assignments.  The positions of the centroids are not arbitrary; they are calculated as the mean of the data points within each cluster.  The context of the course is crucial for understanding the iterative nature of the K-means algorithm and the role of centroids in defining clusters.%
\subsection{Figures}%
\begin{figure}[H]%
\centering%
\end{figure}

%
\newpage%
\section{Slide 30}%
\label{sec:Slide30}%
\subsection{Content}%

\textbf{Python Example}
%
\subsection{Description}%
This slide, part of a lecture on Artificial Intelligence Basics, displays the title "Python Example."  No further content is present on the slide.  The context of the course suggests that this slide is a placeholder for a subsequent slide or slides that will contain Python code examples related to the topic being discussed.  The title indicates that the following content will be a practical demonstration of a concept using the Python programming language.  The lack of any other content implies that the example is yet to be presented.%
\subsection{Figures}%
\begin{figure}[H]%
\centering%
\end{figure}

%
\newpage%
\section{Slide 31}%
\label{sec:Slide31}%
\subsection{Content}%

\textbf{Clustering in Python}
%
\subsection{Description}%
This slide, part of a lecture on Artificial Intelligence Basics, demonstrates K-means clustering using Python.  It shows code snippets and a scatter plot.  The code imports necessary libraries (`matplotlib.pyplot` as `plt` and `seaborn` style, and `make_blobs` from `sklearn.datasets.samples_generator` for creating sample data).  It then generates a dataset (`X`, `y_true`) with 300 samples, 4 centers, and a standard deviation of 0.6, using a random state for reproducibility.  The first scatter plot visualizes the generated data points (`X[:, 0]`, `X[:, 1]`) with a size of 50.  Next, it imports the `KMeans` class from `sklearn.cluster`.  It initializes a `KMeans` object with 4 clusters.  The code fits the model to the data (`X`) using `kmeans.fit(X)`.  Then, it predicts the cluster assignments for each data point (`y_kmeans = kmeans.predict(X)`).  The second scatter plot overlays the data points colored according to their predicted cluster (`y_kmeans`).  Finally, it extracts the cluster centroids (`centers = kmeans.cluster_centers_`) and plots them as larger black points (`plt.scatter(centers[:, 0], centers[:, 1], c='black', s=200, alpha=0.5)`).  The `cmap='viridis'` argument in the second scatter plot specifies a colormap for the clusters. The code snippet demonstrates the application of the K-means algorithm to a dataset and visualizes the resulting clusters and centroids. The context of the course is crucial for understanding the practical implementation of clustering algorithms in Python using the scikit-learn library.%
\subsection{Figures}%
\begin{figure}[H]%
\centering%
\begin{subfigure}[b]{0.15\textwidth}%
\centering%
\includegraphics[width=\textwidth]{/Users/ashoksaravanan/Coding/ScribeLec/server/output/CS 243/lectures/12L2-distance-clusters/figures/31.1.png}%
\caption{}%
\end{subfigure}%
\hspace{0.05\textwidth}%
\begin{subfigure}[b]{0.15\textwidth}%
\centering%
\includegraphics[width=\textwidth]{/Users/ashoksaravanan/Coding/ScribeLec/server/output/CS 243/lectures/12L2-distance-clusters/figures/31.2.png}%
\caption{}%
\end{subfigure}%
\end{figure}

%
\newpage%
\section{Slide 32}%
\label{sec:Slide32}%
\subsection{Content}%

\textbf{How to evaluate Clusters}
%
\subsection{Description}%
This slide, part of a lecture on Artificial Intelligence Basics, presents a title: "How to evaluate Clusters."  It contains no further content.  The title indicates that the following slides will discuss methods for evaluating the quality and effectiveness of clustering results.  The context of the course suggests that this is a transition to a section on evaluating the performance of clustering algorithms, likely after a discussion of clustering methods like K-means.%
\subsection{Figures}%
\begin{figure}[H]%
\centering%
\end{figure}

%
\newpage%
\section{Slide 33}%
\label{sec:Slide33}%
\subsection{Content}%
%
\subsection{Description}%
This slide discusses how to choose the optimal number of clusters (k) in the K-means clustering algorithm.  It begins by explaining the K-means optimization process, which aims to minimize the sum of squared distances from each data point to its assigned cluster centroid. This sum of squared distances is referred to as the "inertia" score, a metric used in the scikit-learn library.  The slide then contrasts a naive approach to choosing k (simply trying different values and selecting the one that minimizes inertia) with a better approach. The naive approach is problematic because the inertia score cannot increase as k increases, leading to the selection of k=n (the total number of data points).  The slide highlights a better approach: selecting the value of k corresponding to the "elbow" point in a plot of inertia versus the number of clusters. This approach balances the accuracy (low inertia) with the complexity (small k) of the clustering solution.  A graph is included, showing a plot of inertia against the number of clusters, illustrating the concept of the "elbow" point. The context of the course is Artificial Intelligence Basics, focusing on practical aspects of machine learning algorithms.%
\subsection{Figures}%
\begin{figure}[H]%
\centering%
\begin{subfigure}[b]{0.15\textwidth}%
\centering%
\includegraphics[width=\textwidth]{/Users/ashoksaravanan/Coding/ScribeLec/server/output/CS 243/lectures/12L2-distance-clusters/figures/33.1.png}%
\caption{}%
\end{subfigure}%
\end{figure}

%
\newpage%
\section{Slide 34}%
\label{sec:Slide34}%
\subsection{Content}%

\textbf{Clustering will return results, no matter the data}
%
\subsection{Description}%
This slide, part of a lecture on Artificial Intelligence Basics, demonstrates the application of different clustering algorithms (K-means, DBSCAN, Complete Link) to a dataset.  Four subplots are shown, each representing a different clustering algorithm applied to the same dataset.  The top-left subplot is labeled "Random Points" and displays a scatter plot of randomly generated data points.  The top-right subplot, labeled "DBSCAN," shows the results of applying the DBSCAN clustering algorithm to the same data.  The bottom-left subplot, labeled "K-means," shows the results of applying the K-means clustering algorithm to the data.  The bottom-right subplot, labeled "Complete Link," shows the results of applying the Complete Link clustering algorithm to the data.  Each subplot shows the data points colored according to the cluster they belong to.  The slide visually illustrates that different clustering algorithms can produce different cluster assignments for the same dataset.  The context of the course is Artificial Intelligence Basics, focusing on the practical application of clustering algorithms and their sensitivity to the choice of algorithm.%
\subsection{Figures}%
\begin{figure}[H]%
\centering%
\begin{subfigure}[b]{0.15\textwidth}%
\centering%
\includegraphics[width=\textwidth]{/Users/ashoksaravanan/Coding/ScribeLec/server/output/CS 243/lectures/12L2-distance-clusters/figures/34.1.png}%
\caption{}%
\end{subfigure}%
\hspace{0.05\textwidth}%
\begin{subfigure}[b]{0.15\textwidth}%
\centering%
\includegraphics[width=\textwidth]{/Users/ashoksaravanan/Coding/ScribeLec/server/output/CS 243/lectures/12L2-distance-clusters/figures/34.2.png}%
\caption{}%
\end{subfigure}%
\hspace{0.05\textwidth}%
\begin{subfigure}[b]{0.15\textwidth}%
\centering%
\includegraphics[width=\textwidth]{/Users/ashoksaravanan/Coding/ScribeLec/server/output/CS 243/lectures/12L2-distance-clusters/figures/34.3.png}%
\caption{}%
\end{subfigure}%
\hspace{0.05\textwidth}%
\begin{subfigure}[b]{0.15\textwidth}%
\centering%
\includegraphics[width=\textwidth]{/Users/ashoksaravanan/Coding/ScribeLec/server/output/CS 243/lectures/12L2-distance-clusters/figures/34.4.png}%
\caption{}%
\end{subfigure}%
\end{figure}

%
\newpage%
\section{Slide 35}%
\label{sec:Slide35}%
\subsection{Content}%

\textbf{Unintuitive k-means results}
%
\subsection{Description}%
This slide, part of a lecture on Artificial Intelligence Basics, displays four scatter plots illustrating the results of applying the K-means clustering algorithm to different datasets.  Each plot demonstrates how the K-means algorithm can produce counterintuitive or unexpected cluster assignments when the data does not meet the algorithm's assumptions.

The top-left plot is labeled "Incorrect Number of Blobs." It shows a dataset with three distinct clusters, but the K-means algorithm, likely with a fixed k=2, has grouped the data into two clusters, misrepresenting the underlying structure.

The top-right plot is labeled "Anisotropy Distributed Blobs."  The data points are clustered along an elongated, non-circular shape.  The K-means algorithm, which assumes spherical clusters, struggles to capture the true cluster structure, resulting in less-than-optimal groupings.

The bottom-left plot is labeled "Unequal Variance."  The data points exhibit clusters with different variances (spread).  The K-means algorithm, which assumes clusters with similar variances, struggles to accurately separate the clusters, leading to overlapping or poorly defined clusters.

The bottom-right plot is labeled "Unevenly Sized Blobs."  The clusters have different sizes.  The K-means algorithm, which treats all clusters equally, may not effectively capture the structure of the data, potentially merging smaller clusters into larger ones or failing to fully separate them.

The slide visually demonstrates the limitations of the K-means algorithm when applied to datasets that do not conform to its underlying assumptions of spherical, isotropic, and equal-sized clusters.  The context of the course is Artificial Intelligence Basics, focusing on the practical application of clustering algorithms and their sensitivity to the characteristics of the data.%
\subsection{Figures}%
\begin{figure}[H]%
\centering%
\begin{subfigure}[b]{0.15\textwidth}%
\centering%

\includegraphics[width=\textwidth]{/Users/ashoksaravanan/Coding/ScribeLec/server/output/CS 243/lectures/12L2-distance-clusters/figures/35.1.png}%
\caption{}%
\end{subfigure}%
\hspace{0.05\textwidth}%
\begin{subfigure}[b]{0.15\textwidth}%
\centering%
\includegraphics[width=\textwidth]{/Users/ashoksaravanan/Coding/ScribeLec/server/output/CS 243/lectures/12L2-distance-clusters/figures/35.2.png}%
\caption{}%
\end{subfigure}%
\hspace{0.05\textwidth}%
\begin{subfigure}[b]{0.15\textwidth}%
\centering%
\includegraphics[width=\textwidth]{/Users/ashoksaravanan/Coding/ScribeLec/server/output/CS 243/lectures/12L2-distance-clusters/figures/35.3.png}%
\caption{}%
\end{subfigure}%
\hspace{0.05\textwidth}%
\begin{subfigure}[b]{0.15\textwidth}%
\centering%
\includegraphics[width=\textwidth]{/Users/ashoksaravanan/Coding/ScribeLec/server/output/CS 243/lectures/12L2-distance-clusters/figures/35.4.png}%
\caption{}%
\end{subfigure}%
\hspace{0.05\textwidth}%
\begin{subfigure}[b]{0.15\textwidth}%
\centering%
\includegraphics[width=\textwidth]{/Users/ashoksaravanan/Coding/ScribeLec/server/output/CS 243/lectures/12L2-distance-clusters/figures/35.5.png}%
\caption{}%
\end{subfigure}%
\end{figure}

%
\newpage%
\section{Slide 36}%
\label{sec:Slide36}%
\subsection{Content}%

\textbf{Interpreting results}
%
\subsection{Description}%
This slide, part of a lecture on Artificial Intelligence Basics, displays a title: "Interpreting results."  No further content is visible on the slide.  The context suggests that the following slides will likely contain information on how to interpret the results of a machine learning algorithm, potentially including visualizations, metrics, and explanations of the findings.  The course is focused on the practical application of AI concepts.%
\subsection{Figures}%
\begin{figure}[H]%
\centering%
\end{figure}

%
\newpage%
\section{Slide 37}%
\label{sec:Slide37}%
\subsection{Content}%

\textbf{Quantitative evaluation}

For prediction tasks there are a variety of quantitative evaluation metrics
\begin{itemize}
    \item Accuracy, squared loss, area under ROC, etc.
\end{itemize}

For cluster analysis there is typically no ground truth
\begin{itemize}
    \item Evaluation aims measure the "goodness" of the resulting clusters
\end{itemize}
%
\subsection{Description}%
This slide discusses quantitative evaluation methods in the context of Artificial Intelligence Basics.  It specifically addresses the differences in evaluation approaches for prediction tasks versus cluster analysis.  For prediction tasks, standard metrics like accuracy, squared loss, and area under the ROC curve are mentioned.  Crucially, for cluster analysis, the slide highlights the absence of a ground truth, unlike prediction tasks.  Evaluation in this case focuses on assessing the "goodness" of the resulting clusters, implying a need for alternative metrics that capture the quality of the clustering structure itself. The context of the course is relevant to understanding the different evaluation strategies for various machine learning tasks.%
\subsection{Figures}%
\begin{figure}[H]%
\centering%
\end{figure}

%
\newpage%
\section{Slide 38}%
\label{sec:Slide38}%
\subsection{Content}%

\textbf{Evaluation Approach: Subjective}

Evaluate how well the clusters "fit" the data
\begin{itemize}
    \item Order the proximity matrix with respect to cluster labels
    \item Inspect visually
\end{itemize}
Good clusterings exhibit clear block pattern
%
\subsection{Description}%
This slide, part of a lecture on Artificial Intelligence Basics, discusses a subjective approach to evaluating the quality of clustering results.  It states that the evaluation focuses on how well the clusters "fit" the data.  Two key steps are outlined: 1) ordering the proximity matrix with respect to cluster labels, and 2) visual inspection.  The slide also notes that good clusterings will exhibit a clear block pattern in the proximity matrix when ordered by cluster labels.  The context of the course is relevant to understanding different evaluation methods for clustering algorithms, which are often qualitative and require human judgment.%
\subsection{Figures}%
\begin{figure}[H]%
\centering%
\end{figure}

%
\newpage%
\section{Slide 39}%
\label{sec:Slide39}%
\subsection{Content}%
%
\subsection{Description}%
This slide, part of a lecture on Artificial Intelligence Basics, presents an example of a "good clustering."  It displays a scatter plot (left) showing data points grouped into three clusters (red, green, and blue).  The x and y axes represent the data features.  The plot visually demonstrates the separation of the data points into distinct clusters.  Accompanying the scatter plot is a heatmap (right) representing a proximity matrix.  The proximity matrix is reordered to reflect the cluster assignments.  The color intensity in the heatmap corresponds to the proximity (similarity) between data points.  The reordering of the matrix is crucial; the blocks of similar colors (e.g., orange, blue) in the heatmap visually demonstrate that data points within the same cluster have high proximity, while data points from different clusters have lower proximity.  The example illustrates a good clustering because the proximity matrix clearly shows distinct blocks corresponding to the clusters in the scatter plot.  The context of the lecture is on evaluating clustering algorithms, where the visual representation of the proximity matrix is used to assess the quality of the clustering results.%
\subsection{Figures}%
\begin{figure}[H]%
\centering%
\begin{subfigure}[b]{0.15\textwidth}%
\centering%
\includegraphics[width=\textwidth]{/Users/ashoksaravanan/Coding/ScribeLec/server/output/CS 243/lectures/12L2-distance-clusters/figures/39.1.png}%
\caption{}%
\end{subfigure}%
\hspace{0.05\textwidth}%
\begin{subfigure}[b]{0.15\textwidth}%
\centering%
\includegraphics[width=\textwidth]{/Users/ashoksaravanan/Coding/ScribeLec/server/output/CS 243/lectures/12L2-distance-clusters/figures/39.2.png}%
\caption{}%
\end{subfigure}%
\hspace{0.05\textwidth}%
\begin{subfigure}[b]{0.15\textwidth}%
\centering%
\includegraphics[width=\textwidth]{/Users/ashoksaravanan/Coding/ScribeLec/server/output/CS 243/lectures/12L2-distance-clusters/figures/39.3.png}%
\caption{}%
\end{subfigure}%
\hspace{0.05\textwidth}%
\begin{subfigure}[b]{0.15\textwidth}%
\centering%
\includegraphics[width=\textwidth]{/Users/ashoksaravanan/Coding/ScribeLec/server/output/CS 243/lectures/12L2-distance-clusters/figures/39.4.png}%
\caption{}%
\end{subfigure}%
\end{figure}

%
\newpage%
\section{Slide 40}%
\label{sec:Slide40}%
\subsection{Content}%
%
\subsection{Description}%
This slide, part of a lecture on Artificial Intelligence Basics, presents an example of a "poor clustering."  It shows a scatter plot on the left, depicting data points colored in red, green, and blue.  The x and y axes represent data features.  The data points are not clearly separated into distinct clusters.  The plot on the right is a heatmap representing a proximity matrix.  The proximity matrix is reordered to reflect the cluster assignments.  However, the heatmap does not exhibit clear blocks of similar colors corresponding to the clusters.  Instead, the colors are mixed and blended, indicating that the proximity between data points within the same cluster is not consistently high, and the proximity between data points from different clusters is not consistently low.  This visual representation demonstrates a poor clustering because the proximity matrix does not clearly reflect the separation of the data into distinct groups. The context of the lecture is on evaluating clustering algorithms, where the visual representation of the proximity matrix is used to assess the quality of the clustering results.  The example contrasts with a "good clustering" example shown in a previous slide, highlighting the difference between well-defined and poorly defined clusters.%
\subsection{Figures}%
\begin{figure}[H]%
\centering%
\begin{subfigure}[b]{0.15\textwidth}%
\centering%
\includegraphics[width=\textwidth]{/Users/ashoksaravanan/Coding/ScribeLec/server/output/CS 243/lectures/12L2-distance-clusters/figures/40.1.png}%
\caption{}%
\end{subfigure}%
\hspace{0.05\textwidth}%
\begin{subfigure}[b]{0.15\textwidth}%
\centering%
\includegraphics[width=\textwidth]{/Users/ashoksaravanan/Coding/ScribeLec/server/output/CS 243/lectures/12L2-distance-clusters/figures/40.2.png}%
\caption{}%
\end{subfigure}%
\hspace{0.05\textwidth}%
\begin{subfigure}[b]{0.15\textwidth}%
\centering%

\includegraphics[width=\textwidth]{/Users/ashoksaravanan/Coding/ScribeLec/server/output/CS 243/lectures/12L2-distance-clusters/figures/40.3.png}%
\caption{}%
\end{subfigure}%
\hspace{0.05\textwidth}%
\begin{subfigure}[b]{0.15\textwidth}%
\centering%
\includegraphics[width=\textwidth]{/Users/ashoksaravanan/Coding/ScribeLec/server/output/CS 243/lectures/12L2-distance-clusters/figures/40.4.png}%
\caption{}%
\end{subfigure}%
\end{figure}

%
\newpage%
\section{Slide 41}%
\label{sec:Slide41}%
\subsection{Content}%

\textbf{Evaluation Approach: Subjective}

Visualize clusters:
\begin{itemize}
    \item Use PCA to reduce dimensionality
    Show only 2 or 3 dimensions
    \item Color points by cluster membership
    \item Inspect results
\end{itemize}
%
\subsection{Description}%
This slide, part of a lecture on Artificial Intelligence Basics, discusses a subjective method for evaluating clustering results.  The approach focuses on visualizing the clusters.  The steps are: 1) Use Principal Component Analysis (PCA) to reduce the dimensionality of the data, showing only 2 or 3 dimensions for easier visualization. 2) Color the data points according to their assigned cluster membership.  3) Inspect the resulting visualization to assess the quality of the clustering.  The slide includes a scatter plot, which is likely the output of the PCA and coloring step, showing a large number of data points colored differently, suggesting different cluster assignments. The context of the lecture is on evaluating clustering algorithms, where visual inspection is a common technique for assessing the quality of the results.%
\subsection{Figures}%
\begin{figure}[H]%
\centering%
\begin{subfigure}[b]{0.15\textwidth}%
\centering%

\includegraphics[width=\textwidth]{/Users/ashoksaravanan/Coding/ScribeLec/server/output/CS 243/lectures/12L2-distance-clusters/figures/41.1.png}%
\caption{}%
\end{subfigure}%
\end{figure}

%
\newpage%
\section{Slide 42}%
\label{sec:Slide42}%
\subsection{Content}%

\textbf{Evaluation Approach: Subjective}

Select a small number of "canonical" members
of each cluster for inspection

\begin{itemize}
    \item E.g., in k-means select the examples that are
    closest to cluster centers
\end{itemize}

Hopefully, the chosen examples can be used to
interpret the characteristics of each cluster (with
domain knowledge)
%
\subsection{Description}%
This slide, part of a lecture on Artificial Intelligence Basics, discusses a subjective approach to evaluating clustering results.  It suggests selecting a small number of representative data points ("canonical" members) from each cluster for closer inspection.  As an example, it mentions using k-means clustering and selecting the data points closest to the cluster centers.  The goal is to use these representative examples to understand the characteristics of each cluster, leveraging domain knowledge to interpret the results. The context of the lecture is on evaluating clustering algorithms, where visual inspection and selection of representative data points are common techniques for assessing the quality and meaning of the clusters.%
\subsection{Figures}%
\begin{figure}[H]%
\centering%
\begin{subfigure}[b]{0.15\textwidth}%
\centering%
\includegraphics[width=\textwidth]{/Users/ashoksaravanan/Coding/ScribeLec/server/output/CS 243/lectures/12L2-distance-clusters/figures/42.1.png}%
\caption{}%
\end{subfigure}%
\end{figure}

%
\newpage%
\section{Slide 43}%
\label{sec:Slide43}%
\subsection{Content}%

\textbf{Evaluation Approach: Objective}

Evaluate the clusters using known class labels

If you have class labels why cluster?
\begin{itemize}
    \item The labeled dataset is usually hand-labeled specifically for evaluation
    \item Cluster the larger dataset automatically
    \item Assess the clustering of the hand labeled dataset in the automatic clustering
\end{itemize}
%
\subsection{Description}%
This slide, part of a lecture on Artificial Intelligence Basics, discusses an objective approach to evaluating clustering results.  It focuses on using known class labels to evaluate the quality of the clustering.  The slide states that if class labels are available, the reason for clustering a dataset is to evaluate the performance of an automatic clustering algorithm on a larger dataset.  The labeled dataset is typically hand-labeled specifically for evaluation purposes.  The steps outlined are: 1) cluster the larger dataset automatically, and 2) assess the clustering of the hand-labeled dataset within the automatically clustered dataset.  The context of the lecture is on evaluating clustering algorithms, where objective evaluation using known class labels is a standard approach.%
\subsection{Figures}%
\begin{figure}[H]%
\centering%
\end{figure}

%
\newpage%
\section{Slide 44}%
\label{sec:Slide44}%
\subsection{Content}%

\textbf{Evaluation via Class Labels: Purity}

Purity: a measure of the extent to which cluster/groups (G) contain objects of a particular class (C)

\begin{itemize}
    \item Each cluster is assigned the class according to the most numerous class of its elements
    \item The purity of each cluster $i$ in $G$ is determined by the number of examples of class $j$ ($N_j$) inside each cluster:
    \[
    \text{purity}(C, G) = \frac{1}{N} \sum_{i=1}^{K} \max_{j} N_j, \quad N_j \in G_i
    \]
\end{itemize}

High purity is easy to achieve when the number of clusters is large

\begin{center}
\begin{tabular}{ccc}
cluster 1 & cluster 2 & cluster 3 \\
\includegraphics[width=2cm]{cluster1.png} & \includegraphics[width=2cm]{cluster2.png} & \includegraphics[width=2cm]{cluster3.png} \\
\end{tabular}
\end{center}

$\frac{1}{17} \cdot (5 + 4 + 3) = 0.71...$
%
\subsection{Description}%
This slide, part of a lecture on Artificial Intelligence Basics, discusses evaluating clustering results using the concept of "Purity."  Purity is a measure of how well the clusters (denoted by $G$) separate objects of different classes (denoted by $C$).  The slide explains that each cluster is assigned the class corresponding to the most frequent class within that cluster.  The purity of a cluster is then calculated as the ratio of the number of examples of the most frequent class within that cluster to the total number of examples in that cluster.  The formula for purity is given as:  $\text{purity}(C, G) = \frac{1}{N} \sum_{i=1}^{K} \max_{j} N_j, \quad N_j \in G_i$.  Here, $N$ is the total number of data points, $K$ is the number of clusters, and $N_j$ is the number of data points belonging to class $j$ within cluster $i$.  The slide emphasizes that high purity is more easily achieved when the number of clusters is large.  The slide includes visual representations of three clusters (cluster 1, cluster 2, and cluster 3) with examples of different classes (represented by 'X' and 'O').  The final calculation shows an example of how purity is computed, using the counts of each class within each cluster.  The context of the lecture is on evaluating clustering algorithms, where purity is a common metric for assessing the quality of the clustering results.%
\subsection{Figures}%
\begin{figure}[H]%
\centering%
\begin{subfigure}[b]{0.15\textwidth}%
\centering%
\includegraphics[width=\textwidth]{/Users/ashoksaravanan/Coding/ScribeLec/server/output/CS 243/lectures/12L2-distance-clusters/figures/44.1.png}%
\caption{}%
\end{subfigure}%
\end{figure}

%
\newpage%
\section{Slide 45}%
\label{sec:Slide45}%
\subsection{Content}%

\textbf{Evaluation via Class Labels: Similarity}

Based on premise that any pair of objects in the same
cluster should have the same class and vice versa

Construct the "observed" similarity matrix

\begin{itemize}
    \item Based on cluster membership
    \item Entry $i, j$ is 1 if $i$ and $j$ are in the same cluster, 0 otherwise
\end{itemize}

Construct the "ideal" similarity matrix

\begin{itemize}
    \item Based on class values
    \item Entry $i, j$ is 1 if $i$ and $j$ are in the same class, 0 otherwise
\end{itemize}

Measure correlation between these two matrices

\begin{tabular}{cccc}
& Point 1 & Point 2 & Point 3 \\
\textbf{Observed} & & & \\
Point 1 & 1 & 0 & 1 \\
Point 2 & 0 & 1 & 0 \\
Point 3 & 1 & 0 & 1 \\
\textbf{Ideal} & & & \\
Point 1 & 1 & 1 & 0 \\
Point 2 & 1 & 1 & 0 \\
Point 3 & 0 & 0 & 1 \\
\end{tabular}
%
\subsection{Description}%
This slide, part of a lecture on Artificial Intelligence Basics, discusses evaluating clustering algorithms using similarity measures based on class labels.  It describes a method for constructing two similarity matrices: an "observed" matrix and an "ideal" matrix.  The "observed" matrix is based on cluster membership, where an entry $i, j$ is 1 if data points $i$ and $j$ are in the same cluster, and 0 otherwise. The "ideal" matrix is based on class labels, where an entry $i, j$ is 1 if data points $i$ and $j$ belong to the same class, and 0 otherwise.  The slide then proposes measuring the correlation between these two matrices as a way to evaluate the quality of the clustering.  The slide includes example matrices, showing the entries for the "observed" and "ideal" similarity matrices.  The context of the lecture is on evaluating clustering algorithms, where this method is used to assess how well the clusters reflect the underlying class structure of the data.%
\subsection{Figures}%
\begin{figure}[H]%
\centering%
\begin{subfigure}[b]{0.15\textwidth}%
\centering%
\includegraphics[width=\textwidth]{/Users/ashoksaravanan/Coding/ScribeLec/server/output/CS 243/lectures/12L2-distance-clusters/figures/45.1.png}%
\caption{}%
\end{subfigure}%
\hspace{0.05\textwidth}%
\begin{subfigure}[b]{0.15\textwidth}%
\centering%
\includegraphics[width=\textwidth]{/Users/ashoksaravanan/Coding/ScribeLec/server/output/CS 243/lectures/12L2-distance-clusters/figures/45.2.png}%
\caption{}%
\end{subfigure}%
\end{figure}

%
\newpage%
\end{document}